\section{Spectral and Multi-resolution Representations}

Fourier methods are natural for periodic flows and turbulence analysis because they describe a field in terms of global frequency modes. They are particularly useful for energy-spectrum analysis and pseudo-spectral solvers. However, global basis functions are less adapted to localized structures and sharp gradients.

Wavelet methods provide a complementary representation localized in both space and scale. This property is useful for RTI fields, where large coherent structures coexist with localized small-scale details. The multi-resolution structure also makes wavelets attractive for compression, denoising, and scale-aware learning.

Wavelet representations are especially relevant for this project because they
combine spatial localization with scale localization. Mallat's work presents
wavelet decompositions as a multi-resolution signal representation able to
separate coarse structures from localized details \citep{Mallat2009}. This
property is directly connected to RTI fields, where large interface
deformations coexist with sharper gradients and smaller mixing structures. In
contrast with purely Fourier representations, wavelets therefore offer a
natural way to introduce scale-aware information into the generative model.

